\chapter{Финансовый аудит на основе размеченных транзакций}

\section{От выписки к управленческому cash flow}

После того как банковская выписка размечена по категориям, можно построить то, чего у компании раньше не было, -- управленческий отчёт о движении денежных средств (ДДС). Принципиальное отличие управленческого ДДС от бухгалтерского: здесь считаются реальные деньги на счёте, а не начисленные доходы и расходы.

Для построения cash flow операции разбиваются на четыре группы:
\begin{itemize}
  \item \textbf{Операционные} -- всё, что связано с основной деятельностью: поступления от клиентов, закупки сырья и упаковки, аренда, персонал, налоги, сервисы.
  \item \textbf{Финансовые} -- кредиты, займы, погашение долга и процентов.
  \item \textbf{Возвраты} -- возвраты от покупателей и возвраты поставщикам.
  \item \textbf{Внутренние переводы} -- перемещения денег между своими счетами.
\end{itemize}

Внутренние переводы обязательно исключаются из операционного денежного потока: если учесть их как расходы и доходы одновременно, картина раздуется в два раза. Это типичная ошибка, которую допускают при первичной обработке банковских выписок без разметки.

\section{Главный финансовый результат}

Результат аудита оказался неожиданным даже в контексте знания о проблемах бизнеса. За 12 месяцев (май~2025 -- апрель~2026):

\begin{center}
\begin{tabular}{lr}
\toprule
\textbf{Показатель} & \textbf{Значение} \\
\midrule
Банковская выручка клиентов & 42~002~899 руб. \\
Операционные списания & 41~980~081 руб. \\
\textbf{Operating cash flow} & \textbf{22~818 руб.} \\
Денежная операционная маржа & \textbf{0,05\%} \\
\bottomrule
\end{tabular}
\end{center}

Сорок два миллиона рублей оборота -- и двадцать две тысячи рублей на выходе. Это не убыток и не банкротство, но это и не бизнес, который формирует доход для собственников. Денежная маржа в 0,05\% означает, что компания буквально балансирует на грани: каждый крупный незапланированный платёж или задержка от клиента моментально уводит денежный поток в минус.

\section{Структура расходов}

Чтобы понять, куда уходят деньги, посмотрим на структуру операционных списаний. Ключевые статьи:

\begin{center}
\begin{tabular}{lrr}
\toprule
\textbf{Категория расходов} & \textbf{Сумма, руб.} & \textbf{Доля от выручки} \\
\midrule
Производственные закупки & 29~160~903 & 69,4\% \\
Упаковка & 3~222~206 & 7,7\% \\
\textbf{Итого прямые закупки} & \textbf{32~383~108} & \textbf{77,1\%} \\
\addlinespace
Аренда и коммунальные & — & — \\
Персонал & — & — \\
Налоги и взносы & — & — \\
Сервисы и прочее & — & — \\
\bottomrule
\end{tabular}
\end{center}

Главная статья -- \textbf{производственные закупки}: почти 70 копеек с каждого рубля выручки уходит на сырьё и ингредиенты. Это не само по себе плохо для производственного бизнеса, но означает, что любое повышение цен на сырьё немедленно бьёт по марже.

Прокси-маржа после прямых денежных закупок составляет:
\[
\text{Прокси-маржа} = \frac{42\,002\,899 - 32\,383\,108}{42\,002\,899} \approx 22{,}9\%
\]

22,9\% -- это не бухгалтерская валовая прибыль (она не учитывает складские остатки, незавершённое производство и начисленную амортизацию). Это приближённый денежный вклад, который остаётся после оплаты сырья и упаковки. Именно этот вклад должен покрыть аренду, персонал, налоги и всё остальное. Как показывает итоговый operating cash flow, он едва справляется с этой задачей.

\section{Кассовые провалы по месяцам}

Помесячный анализ денежного потока выявил шесть месяцев, когда операционный денежный поток уходил в минус:

\begin{center}
\begin{tabular}{lc}
\toprule
\textbf{Месяц} & \textbf{Знак OCF} \\
\midrule
Май 2025 & отрицательный \\
Июль 2025 & отрицательный \\
Сентябрь 2025 & отрицательный \\
Декабрь 2025 & отрицательный \\
Февраль 2026 & отрицательный \\
Апрель 2026 & отрицательный \\
\bottomrule
\end{tabular}
\end{center}

Интересно, что провалы происходят примерно раз в два месяца. Это может указывать на структурный паттерн: циклические закупочные всплески, сезонность клиентских оплат или задержки расчётов с конкретными сетями. Для планирования важно, что эти месяцы заранее предсказуемы -- а значит, к ним можно готовиться: формировать резерв ликвидности или переносить часть закупок.

\section{Анализ концентрации контрагентов}

Концентрация в финансовом анализе оценивается через \textbf{индекс Херфиндаля--Хиршмана (HHI)} и метрики CR:

\[
\text{HHI} = \sum_{i=1}^{n} s_i^2, \quad
\text{CR}_k = \sum_{i=1}^{k} s_i
\]

где $s_i$ -- доля $i$-го контрагента в суммарном потоке. HHI в пунктах (умноженный на 10~000): ниже 1~500 -- низкая концентрация, 1~500--2~500 -- высокая, выше 2~500 -- очень высокая.

\subsection{Концентрация расходов}

По поставщикам и платёжным контрагентам компания сильно зависит от нескольких ключевых партнёров:

\begin{center}
\begin{tabular}{lr}
\toprule
\textbf{Метрика} & \textbf{Значение} \\
\midrule
HHI (поставщики) & 1~541 пунктов \\
CR1 & 34,1\% \\
CR5 & 68,0\% \\
CR10 & 84,7\% \\
\addlinespace
Крупнейший поставщик (ООО <<РАДУГА~Н>>) & 14~387~180 руб. \\
\bottomrule
\end{tabular}
\end{center}

HHI на уровне 1~541 -- это высокая концентрация. Один поставщик занимает треть всех дебетовых платежей. Если он повысит цены, изменит условия отсрочки или просто не сможет вовремя отгрузить сырьё -- это немедленно скажется на производстве и кассе.

\subsection{Концентрация поступлений}

Ситуация с покупателями интересна тем, что два уровня анализа дают разные ответы.

На уровне буквальных банковских названий контрагентов концентрация выглядит \textit{низкой}: HHI всего 82 пункта. Но это артефакт структуры сети: один крупный покупатель -- торговая сеть <<Светофор>> -- разбит на множество юридических лиц и расчётных счетов. В банковской выписке они выглядят как разные контрагенты.

После агрегации по данным 1С и группировки по брендам картина меняется кардинально:

\begin{center}
\begin{tabular}{lr}
\toprule
\textbf{Уровень анализа} & \textbf{HHI, пунктов} \\
\midrule
Банковская выписка (буквальные названия) & 82 \\
1С (юридические лица) & 2~591 \\
По брендам/сетям & 7~444 \\
\addlinespace
Доля бренда <<Светофор>> & 85,3\% оплат \\
\bottomrule
\end{tabular}
\end{center}

HHI по брендам 7~444 -- это запредельная концентрация. По сути, судьба денежного потока компании почти полностью зависит от одной торговой сети. Это не проблема сегодняшнего дня, но это серьёзный стратегический риск: изменение условий работы со стороны сети, задержка оплат на месяц или пересмотр ассортиментной матрицы могут обрушить кассу.

\section{Сверка банка и 1С}

Сверка данных из банковской выписки с оплатами в системе 1С -- обязательный шаг для подтверждения корректности данных. Результаты:

\begin{center}
\begin{tabular}{lr}
\toprule
\textbf{Источник} & \textbf{Сумма, руб.} \\
\midrule
Кредитовые поступления банка & 42~010~518 \\
1С-оплаты & 42~002~971 \\
Разница & $-$7~547 \\
\bottomrule
\end{tabular}
\end{center}

На уровне итоговых сумм за год банк и 1С практически совпадают: расхождение менее восьми тысяч рублей на сорок два миллиона оборота -- это 0,018\%. По документам совпадают 2~621 позиция; суммарное отклонение по сопоставленным документам составляет всего 72 рубля, максимальное расхождение по одной строке -- 1,13 руб.

Основная причина разладки технического характера: 1С хранит суммы оплат в округлённых рублях, банк содержит копейки. Кроме того, названия контрагентов в банковской выписке включают расчётные счета и наименование банка, а в 1С -- иногда адрес торговой точки. Поэтому автоматическая построчная сверка без нормализатора контрагентов будет давать ложные расхождения там, где их нет.

\section{Выводы финансового аудита}

Финансовый аудит позволяет сформулировать несколько чётких диагнозов.

Во-первых, \textbf{проблема компании -- не отсутствие выручки, а сверхтонкая денежная маржа}. Оборот есть, но от сорока двух миллионов рублей поступлений остаётся около двадцати трёх тысяч рублей операционного денежного потока. Это нулевая маржа безопасности.

Во-вторых, \textbf{главная зона оптимизации -- прямые закупки}. Производственное сырьё и упаковка занимают вместе 77\% выручки. Снижение закупочных цен даже на несколько процентов дало бы заметный прирост кассового потока.

В-третьих, \textbf{кассовые провалы цикличны и предсказуемы}. Шесть месяцев с отрицательным OCF равномерно распределены по году, что указывает на структурный паттерн. Это открывает возможность для профилактики: платёжный календарь, резерв ликвидности, заблаговременный перенос части закупок.

В-четвёртых, \textbf{зависимость от одного покупателя -- серьёзный риск}. Восемьдесят пять процентов оплат от одной сети -- это не диверсификация, а уязвимость. Расширение клиентской базы -- стратегический приоритет.
